\documentclass[11pt, a4paper]{article}

\usepackage[T1]{fontenc}
\usepackage[utf8]{inputenc}
\usepackage{tgpagella}
\usepackage[spanish, es-noshorthands]{babel}
\usepackage{amsmath, amssymb}
\usepackage[most]{tcolorbox}
\usepackage[margin=2.5cm]{geometry}
\usepackage{fancyhdr}

\pagestyle{fancy}
\fancyhf{}
\setlength{\headheight}{16pt}
\lhead{\textbf{UCB Sede Tarija} | Ing. Mecatrónica}
\rhead{IMT-221 | Soluciones S06-3}
\renewcommand{\headrulewidth}{0.4pt}

\begin{document}

\section*{Solucionario y Guía Docente: Interacción Multietapa}

\subsection*{Ejercicio Diagnóstico}
1. $g_m = 1\,\text{mA} / 25.85\,\text{mV} \approx 38.68\,\text{mS}$. \\
2. $r_\pi = 100 / 0.03868 \approx 2.585\,\text{k}\Omega$. \\
3. \textit{[Revisar dibujo del estudiante: B a E ($r_\pi$), C a E (fuente $g_m v_\pi$ apuntando abajo), E en Tierra AC].} \\
4. $A_{v0} = -(38.68\text{mS})(4.7\text{k}\Omega) \approx -181.8$. \\
5. Suposición: Emisor en tierra AC, $r_o \to \infty$, señal en la base ideal (sin resistencia de fuente).

\subsection*{Ejercicio Integrado Fuente-Amplificador-Carga}
1. $g_m \approx 38.68\,\text{mS}$ y $r_\pi \approx 2.585\,\text{k}\Omega$. \\
2. $R_{in} = R_B \parallel r_\pi = 100\text{k}\Omega \parallel 2.585\text{k}\Omega \approx 2.520\,\text{k}\Omega$. \\
3. $v_i$ será menor debido a la atenuación de $R_S$. $v_i/v_s = 2.520 / (10 + 2.520) \approx 0.2013$. \\
\textbf{Error común:} Asumir $v_s = v_i$. \textit{Sonda pedagógica:} "¿Qué resistencia existe entre el generador ideal y el nodo de entrada del amplificador?" \\
4. $R_L$ reducirá la magnitud de la ganancia. $R_C \parallel R_L = 4.7\text{k} \parallel 10\text{k} \approx 3.197\,\text{k}\Omega$. \\
\textbf{Error común:} Usar $-g_m R_C$ ignorando $R_L$. \textit{Sonda:} "¿Qué resistencia paralela real ve la fuente dependiente en el colector?" \\
5. $A_v = -g_m (R_C \parallel R_L) = -(0.03868)(3197) \approx -123.7$. \\
6. $G_v = (v_i/v_s) A_v = (0.2013)(-123.7) \approx -24.9$. \\
7. Tendería a $A_{v0} = -181.8$. \\
8. Si $R_S \to 0$, entonces $v_i \approx v_s$, y $G_v$ subiría, igualando la magnitud de $A_v$ cargada (-123.7).

\subsection*{Preguntas Conceptuales de Carga y Topología}
\textbf{Selección de Topologías:}
\begin{itemize}
    \item \textbf{Buffer entre fuente y carga pesada:} Colector Común (CC). \textit{Sonda:} "¿Qué sucede con la corriente extraída de la fuente si usamos un CE directo?"
    \item \textbf{Por qué no multiplicar ganancias descargadas en cascada:} La Etapa 2 modifica la resistencia de carga efectiva vista por la Etapa 1. Multiplicar $A_{v0,1} \times A_{v0,2}$ produce un error gigante por exceso.
\end{itemize}

\end{document}
